\documentclass[11pt]{article}

\usepackage[a4paper,margin=25mm]{geometry}
\usepackage[T1]{fontenc}
\usepackage{lmodern}
\usepackage{amsmath,amssymb,amsthm,mathtools}
\usepackage{booktabs,array,longtable,pdflscape,capt-of}
\usepackage{enumitem}
\usepackage{microtype}
\usepackage{xcolor}
\usepackage[hidelinks]{hyperref}

\definecolor{proofblue}{RGB}{26,55,92}
\hypersetup{
  pdftitle={The Primitive Generalized Fermat Equation of signature 3,5,7: From the published frontier to a certified proof},
  pdfsubject={A computer-assisted proof},
  colorlinks=true,
  linkcolor=proofblue,
  citecolor=proofblue,
  urlcolor=proofblue
}

\newcommand{\Q}{\mathbf Q}
\newcommand{\Z}{\mathbf Z}
\newcommand{\F}{\mathbf F}
\newcommand{\Pone}{\mathbf P^1}
\newcommand{\OO}{\mathcal O}
\newcommand{\m}{\mathfrak m}
\newcommand{\pp}{\mathfrak p}
\newcommand{\PP}{\mathfrak P}
\newcommand{\Sha}{\mathop{\mathrm{Sha}}}
\newcommand{\rank}{\operatorname{rank}}
\newcommand{\ord}{\operatorname{ord}}
\newcommand{\Gal}{\operatorname{Gal}}
\newcommand{\disc}{\operatorname{disc}}
\newcommand{\Res}{\operatorname{Res}}

\newtheorem{theorem}{Theorem}[section]
\newtheorem{proposition}[theorem]{Proposition}
\newtheorem{lemma}[theorem]{Lemma}
\newtheorem{corollary}[theorem]{Corollary}
\theoremstyle{definition}
\newtheorem{definition}[theorem]{Definition}
\theoremstyle{remark}
\newtheorem{remark}[theorem]{Remark}

\title{\textbf{The Primitive Generalized Fermat Equation}\\[1mm]
       $x^3+y^5=z^7$\\[3mm]
       \large From the published frontier to a certified proof}
\author{\mbox{}}
\date{13 September 2026}

\begin{document}
\maketitle

\begin{abstract}
We prove that the generalized Fermat equation
\[
                         x^3+y^5=z^7
\]
has no solution in nonzero coprime integers, and we make explicit how the
proof grows from the previously published frontier.  Dahmen--Siksek had
established the signed local descent and eliminated the three cases in which
the associated degree-seven algebra is reducible.  Putz had then proved that
the remaining irreducible case can involve only seven septic fields: six
pure fields and one exceptional field.  The missing step was to exclude
those seven fields.

For the six pure fields we construct an explicit Fano resolvent, descend to
a smooth plane quartic of genus three over $\Q(\sqrt{-7})$, and prove that
its rational-parameter locus consists only of five points above the branch
values.  For the
exceptional field we combine the modularity and conductor results of
Pacetti--Villagra Torcomian with level lowering over $\Q(\sqrt5)$ and a
finite Hecke comparison at $29$.  We also reconstruct the three published
reducible-sector arguments as independently replayable calculations, adding
a database-free identification of the quadratic field and an intrinsic
formal-group treatment at the ramified prime.  The paper labels every major
step as published input, reconstructed input, or new argument.  All
project-specific finite calculations are supplied as exact certificates,
programs, inputs, and authenticated logs.
\end{abstract}

\tableofcontents

\section{Introduction: the frontier and the completion}

\subsection{The equation and the theorem}

For pairwise coprime nonzero integers, the equation
\begin{equation}\label{eq:original}
                         x^3+y^5=z^7
\end{equation}
is the generalized Fermat equation of signature $(3,5,7)$.  We shall use
the variable order
\begin{equation}\label{eq:DS-order}
                    (X,Y,Z)=(y,x,z),\qquad X^5+Y^3=Z^7.
\end{equation}
All signs are allowed.  Since a prime dividing two of $X,Y,Z$ also divides
the third, primitivity is equivalent here to pairwise coprimality.  Our main
result is the following.

\begin{theorem}\label{thm:main}
There are no nonzero primitive integers $x,y,z$ satisfying
\[
                             x^3+y^5=z^7.
\]
\end{theorem}

The purpose of this introduction is not only to state the theorem, but to
say precisely what was known before the present work and what had still to
be done.  This distinction matters in a long computer-assisted proof: it
makes clear which general theorems are imported, which earlier computations
are reconstructed for confidence, and which arguments close the remaining
gap.

\subsection{The published frontier}

The mathematical starting point was already far advanced.  Dahmen and
Siksek proved the following strong partial result
\cite[Theorem~1]{DahmenSiksek}: every nontrivial primitive solution of
$X^5+Y^3=Z^7$ satisfies one of the alternatives
\begin{equation}\label{eq:DS-theorem-one}
 2\cdot3\cdot5\mid Z\ \text{ and }\ 7\nmid XY,
 \qquad\text{or}\qquad
 2\nmid Z, 3\nmid XYZ, 5\nmid YZ, 7\mid Y.
\end{equation}
In particular, a primitive solution for which $3\mid XY$, $5\mid Y$, or
$7\mid X$ is trivial.  Their descent introduces the degree-seven algebra
used below, proves that only four factor-degree patterns can occur, and
eliminates the three reducible patterns
\[
                         [1,6],\qquad[2,5],\qquad[3,4].
\]
The only surviving pattern is therefore the irreducible one, $[7]$.

Putz then made this residual problem finite.  His exhaustive Hunter
enumeration, with the prescribed local algebra data, proves that an
irreducible descent algebra must be one of exactly seven fields
\cite[Theorem~3.50]{Putz}.  Six are pure septic fields and the seventh is an
exceptional septic field.  Thus the published frontier can be summarized in
one sentence:
\begin{quote}
\emph{a nontrivial primitive solution would have to pass through one of seven
explicit irreducible degree-seven fields.}
\end{quote}
This is a remarkably narrow frontier, but it is not yet a contradiction.
One must still connect each field to an arithmetic object on which rational
points or residual Galois representations can be exhausted.

\subsection{What is added here}

The principal new work is the elimination of those seven fields.

\begin{enumerate}[label=\textnormal{(\roman*)},leftmargin=2.2em]
\item For the six pure fields we construct the missing geometric bridge.
The thirty Fano planes attached to the septic cover split into two
$A_7$-orbits.  Their compatible pairs give a degree-$120$ resolvent whose
quotient is an explicit smooth plane quartic $C_+$ over
$\Q(\sqrt{-7})$.  A rational-parameter theorem shows that the relevant
parameter on $C_+$ takes rational values only at five points above
$0$, $1$, and $\infty$.  A nonzero primitive solution would have parameter
outside that set, so all six pure fields are excluded.

\item For the exceptional field we place the specialization in the Frey
family of Pacetti--Villagra Torcomian, use their modularity and conductor
theorems together with Hilbert level lowering, and reduce both the
irreducible and reducible residual alternatives to finite lists.  At a prime
above $29$, the four curve trace polynomials and the six possible modular or
character factors have $24$ nonzero resultants modulo $7$.  This excludes the
exceptional field.

\item We reconstruct the previously published reducible sectors in an exact
certificate framework.  This is not presented as a new elimination of those
sectors.  It provides an independent route through the finite arithmetic and
adds two genuine refinements: the quadratic factor is identified from its
intrinsic local labels without a field database, and the ramified $7$-adic
step is carried out inside the formal neighbourhood of the given integral
Jacobian model.

\item We compose all four sectors into a single machine-checkable theorem
graph.  Every project-specific finite assertion is tied to exact input bytes,
an explicit predicate, and a replay program or authenticated completed log.
Published general theorems remain ordinary bibliographic dependencies.
\end{enumerate}

\subsection{A provenance-aware proof map}

The following table is the shortest guide to the logical ownership of the
proof.  ``Reconstructed'' means that this paper supplies an independent
exact certificate for a result whose mathematical argument was already in
the literature; it does not claim priority for that result.

\begin{center}
\small
\renewcommand{\arraystretch}{1.18}
\begin{tabular}{@{}>{\raggedright\arraybackslash}p{.19\textwidth}
                    >{\raggedright\arraybackslash}p{.25\textwidth}
                    >{\raggedright\arraybackslash}p{.46\textwidth}@{}}
\toprule
stage & status before this paper & role in this paper\\
\midrule
signed local descent and four factor sectors
 & published by Dahmen--Siksek
 & imported, restated, and independently checked\\
sectors $[1,6]$, $[2,5]$, $[3,4]$
 & eliminated by Dahmen--Siksek
 & reconstructed with certificates; two local refinements are added\\
irreducible sector $[7]$
 & reduced by Putz to seven explicit fields
 & the seven-field list is imported and authenticated\\
six pure septic fields
 & not eliminated at the published frontier
 & excluded here by the Fano resolvent and the rational-parameter theorem\\
one exceptional septic field
 & modular tools available, but no terminal comparison
 & excluded here by level lowering, ray characters, and the $q=29$ resultant test\\
global theorem
 & one finite irreducible frontier remained
 & all four sectors are composed into the contradiction of Theorem~\ref{thm:main}\\
\bottomrule
\end{tabular}
\normalsize
\end{center}

\subsection{How the proof was completed}

The working principle was to treat a hypothetical solution not as a large
integer triple to be chased, but as one mathematical state seen through
several exact representations.  The same state has a descent algebra, local
factor labels, ideal classes, divisor classes, finite-quotient coordinates,
and $p$-adic analytic coordinates.  Each representation supplies a
constraint.  The proof succeeds when the intersection of all compatible
states is empty.

This viewpoint dictated a calm order of work:
\begin{enumerate}[leftmargin=2.2em]
\item fix the strongest published frontier before beginning a new search;
\item translate each remaining assertion into a finite interface with exact
inputs and a stated acceptance predicate;
\item use exploratory computation only to discover a model, prime, or
quotient, and then replace discovery by a deterministic certificate;
\item keep the four descent sectors separate until each has a terminal
statement, and only then compose them;
\item distinguish a recomputation of a published argument from a genuinely
new bridge or obstruction.
\end{enumerate}
This is why the proof is modular.  A specialist can inspect the new
irreducible argument without repeating the reducible sectors, or can replay
one finite computation without trusting the remainder of the software tree.

The central object throughout is the polynomial
\begin{equation}\label{eq:phi}
 \Phi(T)=15T^7-35T^6+21T^5,
 \qquad \eta=\frac{X^5}{Z^7},
 \qquad f_\eta(T)=\Phi(T)-\eta.
\end{equation}
The associated etale algebra
\[
                         A_\eta=\Q[T]/(f_\eta)
\]
encodes a hypothetical solution in several complementary forms.  Local
factorization describes which global fields may occur.  Descent maps those
fields to curves.  Mordell--Weil information restricts the resulting points
to finitely many congruence classes, and local analytic functions decide
whether those classes contain a point compatible with the original rational
parameter.

The argument uses standard published results in descent, number-field
enumeration, modularity, level lowering, formal groups, Mordell--Weil sieves,
and Chabauty methods.  The large finite calculations specific to
\eqref{eq:original} are exact and are separated into independently
checkable certificates.  Thus the proof has two visible layers: general
theorems, cited in the usual way, and finite arithmetic, supplied with the
paper.  Sections~\ref{sec:descent}--\ref{sec:reducible-end} reconstruct the
descent and reducible sectors.  Section~\ref{sec:irreducible} begins at
Putz's seven-field frontier and contains the new terminal argument.  The
final sections compose the proof and describe the certificates.

\section{The degree-seven descent algebra}\label{sec:descent}

\paragraph{Provenance.}
The local covering and exhaustive four-sector reduction in this section are
published inputs from Dahmen--Siksek.  We record them because they are the
interface between a primitive solution and every later computation.

The derivative and the two branch values are given by
\begin{equation}\label{eq:derivative}
          \Phi'(T)=105T^4(T-1)^2,
          \qquad \Phi(0)=0,\quad \Phi(1)=1.
\end{equation}
In particular, $\Phi$ is strictly increasing on $\mathbb R$.  For a
nonzero solution, $\eta\notin\{0,1\}$, and $f_\eta$ is separable with one
real root.  A resultant computation gives
\begin{equation}\label{eq:discriminant}
 \disc(f_\eta)
   =-3^6 5^6 7^7\eta^4(\eta-1)^2.
\end{equation}
Its squareclass is therefore $-7$.

The local covering theorem at $7$ in \cite[Proposition~3.5]{DahmenSiksek}
gives local factor-degree patterns $[1,6]$, $[2,5]$, $[3,4]$, and $[7]$.
Global factors induce groupings of the irreducible local factors, so their
degrees are obtained by coarsening these patterns.  Consequently the
possible global factor-degree patterns are exactly among
\begin{equation}\label{eq:four-sectors}
                  [1,6],\qquad [2,5],\qquad [3,4],\qquad [7].
\end{equation}
In particular, there are at most two global factors.  For example, if
$7\mid X$ and $r=\ord_7(X)$, the Newton polygon has vertices
$(0,5r),(5,1),(7,0)$, giving irreducible local factors of degrees $5$ and $2$.
If $7\mid Y$ and $r=\ord_7(Y)$, translation by $1$ gives vertices
$(0,3r),(3,1),(7,0)$, giving irreducible local factors of degrees $3$ and $4$.
The unit case, and the $7\mid Z$ case after scaling to an integral model,
give the remaining local patterns by the cited proposition.

We now eliminate the four sectors one at a time.

\section{The rational-factor sector}

\paragraph{Provenance.}
The mathematical elimination in this section is due to Dahmen--Siksek.  This
section gives a certificate-level reconstruction of their
Proposition~6.1 and Lemmas~6.2--6.8 \cite[Section~6]{DahmenSiksek}.

Assume that $r\in\Q$ is a root of $f_\eta$, and put $w=rZ$.
Clearing denominators gives
\begin{equation}\label{eq:rational-root-cleared}
       15w^7-35Zw^6+21Z^2w^5-X^5=0.
\end{equation}
The translated identity
\[
 \Phi(T)-1=(T-1)^3(15T^4+10T^3+6T^2+3T+1)
\]
also gives
\begin{equation}\label{eq:quartic-cube}
 (w-Z)^3
 (Z^4+3Z^3w+6Z^2w^2+10Zw^3+15w^4)=-Y^3.
\end{equation}

\subsection{Local restrictions and quadratic ideal descent}

Newton polygons at $3$, $5$, and $7$ give
\begin{equation}\label{eq:rational-local}
 3\nmid Z,\qquad 5\nmid YZ,\qquad 7\nmid Y,\qquad
 \ord_3(r)\in\{0,1,-1\},\qquad \ord_5(r)\geq0.
\end{equation}
The rational-root theorem applied to \eqref{eq:rational-root-cleared}
then shows that $w$ is an integer or one third of an integer.  Moreover, for
every prime $p$, at least one of $\ord_p(Z)$ and $\ord_p(w)$ is zero.

For all three reducible sectors fix the notation
\begin{equation}\label{eq:k-minus35}
 k=\Q(\omega)=\Q(\sqrt{-35}),\qquad
 \omega^2+\omega+9=0,\qquad \delta=2\omega+1=\sqrt{-35}.
\end{equation}
Thus $\OO_k=\Z[\omega]$.  Put
\[
 A=Z^2+wZ-3w^2,\qquad B=-wZ-w^2,\qquad \mathcal F=A+B\omega.
\]
Then
\[
 \mathcal F\overline{\mathcal F}=A^2-AB+9B^2
        =Z^4+3Z^3w+6Z^2w^2+10Zw^3+15w^4.
\]
The class group of $k$ is $C_2$.  The prime $3$ splits, and exact ideal
arithmetic reduces \eqref{eq:quartic-cube} to three possibilities:
\begin{equation}\label{eq:three-cubic-cases}
 \mathcal F=\alpha(u+v\omega)^3,\qquad
               \alpha\in\{1,9\omega,3\omega^2\},\qquad u,v\in\Q.
\end{equation}
For $\alpha=1$, \cite[Lemma~6.4]{DahmenSiksek} further gives
$w,u,v\in\Z$, with $\gcd(Z,w)=\gcd(u,v)=1$.

Expanding \eqref{eq:three-cubic-cases} produces three genus-two curves.
One has no $\Q_2$-point because its primitive homogeneous sextic is always
$5$ modulo $8$.  A second has precisely the two rational points
$(-2,\pm540)$; its Jacobian has rank one and the rational points are
exhausted by Chabauty's method.  Substitution gives $r=-1/4$ and
\[
                         \Phi(-1/4)=-\frac{491}{2^{14}},
\]
which cannot be a reduced quotient of a fifth power by a seventh power.

\subsection{The main genus-two curve}

The remaining case gives the curve
\begin{equation}\label{eq:rational-genus2}
 C_1:\quad
 y^2=x^6-6x^5+69x^4-38x^3+6x^2-240x+1057.
\end{equation}
Its Jacobian is torsion-free of rank three, and three explicit Mumford
divisors form a $\Z$-basis.  A Mordell--Weil sieve in the sense of
\cite{BruinStollSieve}, modulo $1080J_1(\Q)$,
places every rational point in one of the six classes represented by
\[
 \infty_+,\ \infty_-,\ (1/2,\pm245/8),\ (3,\pm65).
\]
Exact formal-group calculations on these classes imply, for every affine
$(x,y)\in C_1(\Q)$,
\begin{equation}\label{eq:rational-local-alternative}
                         \ord_3(x)\leq-3
                    \quad\hbox{or}\quad
                         \ord_5(y)\geq1.
\end{equation}

The point on \eqref{eq:rational-genus2} arising from
\eqref{eq:three-cubic-cases} has
\[
 x=\frac uv,\qquad
 y=\frac{(Z+w)^2+3w^2}{v^3}.
\]
The case $v=0$ is also visible on the projective model.  Coprimality forces
$u=\pm1$, while the two descent equations give
\[
       Z^2+wZ-3w^2=\pm1,\qquad -w(Z+w)=0.
\]
If $w=0$, then $X=0$; if $Z=-w$, then $-3w^2=\pm1$.  Both are impossible
under the hypotheses.

For $v\neq0$, the second arm of
\eqref{eq:rational-local-alternative} would imply
$5\mid (Z+w)^2+3w^2$, impossible for a primitive pair modulo $5$.
Hence $27\mid v$.  The second descent equation gives
$27\mid w(Z+w)$.  Coprimality and \eqref{eq:rational-local} exclude
$27\mid w$, so $27\mid Z+w$.  Reducing the first descent equation modulo
$3$ now gives $3\mid u$, contradicting $\gcd(u,v)=1$.

\begin{proposition}\label{prop:rational-sector}
The polynomial $f_\eta$ has no rational root.  Hence the sector $[1,6]$
is empty.
\end{proposition}

\section{The quadratic-factor sector}

\paragraph{Provenance.}
The reduction to the fibre product and its six rational points is
Dahmen--Siksek, Proposition~5.1 and Lemma~5.2
\cite[Section~5]{DahmenSiksek}.  We include the construction because it is
the complete bridge from a quadratic root to the final fifth-power test.  The
database-free recovery of $\Q(\sqrt{-35})$ from its local labels is the
additional refinement used here.

Suppose that $A_\eta$ has a quadratic direct factor $L$.  Its real
signature, ramification, and dyadic labels identify it without a field
database.

Since $f_\eta$ has exactly one real root, $L$ is imaginary.  At every prime
outside $\{3,5,7\}$, the scaled local models are separable and every direct
factor is unramified.  At $2$, the possible local degree patterns are
$[1,2,4]$ and $[1,3,3]$.  A quadratic global direct factor must therefore
be the unique unramified quadratic component in the first pattern.  Thus
$2$ is inert in $L$.  The negative odd fundamental discriminants supported
on $\{3,5,7\}$ and inert at $2$ are $-3$ and $-35$.  Finally, the
degree-two Newton-polygon component at $7$ is ramified.  Therefore
\begin{equation}\label{eq:quadratic-field}
                              L=k=\Q(\omega).
\end{equation}

Let $u$ be a nonrational root of $f_\eta$ in $L$ and $v=\bar u$.  Symmetric
elimination gives a rational point on
\begin{equation}\label{eq:fibre-product}
 D:\quad
 \begin{cases}
 \beta^2=\alpha^4-4\alpha^3+6\alpha^2+9,\\
 \alpha(\alpha-4)=-35\gamma^2.
 \end{cases}
\end{equation}
Here
\[
                    \alpha=\frac{(u+v)^2}{uv}.
\]

\begin{proposition}\label{prop:six-points}
The rational points on $D$ are exactly
\[
 (0,\pm3,0),\qquad
 \left(\frac{35}{9},\ \pm\frac{782}{81},\ \pm\frac19\right),
\]
with independent signs in the second family.
\end{proposition}

\begin{proof}
Map $D$ to the genus-two curve
\begin{equation}\label{eq:quadratic-genus2}
 -35y^2=x(x-4)(x^4-4x^3+6x^2+9).
\end{equation}
The etale algebra of the right side is
$\Q\times\Q\times L_4$, where
\[
 L_4=\Q(\vartheta),\qquad
 \vartheta^4-4\vartheta^3+6\vartheta^2+9=0.
\]
An exact $2$-cover descent \cite{BruinStollDescent} gives three
fake-Selmer classes, whose components
at $0$, $4$, and $\vartheta$ are
\[
\begin{array}{c|ccc}
 &0&4&\vartheta\\ \hline
 \xi_1&-3&1&3(\vartheta-4)\\
 \xi_2&-35&1&1\\
 \xi_3&-35&1&-5.
\end{array}
\]
The first class cannot lift to a nonendpoint rational point on $D$.
Indeed, in this class $x=-3ta^2$ and $x-4=tb^2$, whereas the additional
condition on $D$ requires $x(x-4)/(-35)$ to be a square.  This would make
$3/35$ a square, which is impossible already over $\Q_7$.  The endpoint
$x=4$ does not lift either, since it would require $\beta^2=105$.
The other two classes map to elliptic curves over $L_4$.

Their
ranks are respectively $3$ and $2$, their torsion is trivial, and elliptic
Chabauty \cite{BruinEC} retains only the rational $x$-coordinates
\[
                               x=0,\qquad x=35/9.
\]
For the rank-three curve, completeness is certified by saturation at every
prime dividing the Chabauty index
\[
\begin{split}
6592880678883323322765152793600
={}&2^{10}3^5 5^2 7^4\cdot11\cdot17\cdot19\cdot29\cdot31\cdot37\\
 &\cdot43\cdot47\cdot53\cdot61\cdot71\cdot83\cdot97.
\end{split}
\]
Substitution into \eqref{eq:fibre-product} gives precisely the six points
stated.
\end{proof}

If $\alpha=0$, then $v/u=-1$.  The root relation gives either $u=0$ or
$u=\pm\delta/5$.  The nonzero alternatives satisfy
\[
                            \Phi(u)=\frac{7^4}{5^2},
\]
which is not a reduced quotient of a fifth power by a seventh power.  If
$\alpha=35/9$, then
\[
 \frac vu=\frac{17\pm\delta}{18}.
\]
The root relation gives either a branch value or the rational number
\[
 \frac{495535757310206240323260222}
      {432344158742282808792052775},
\]
whose coprime numerator and denominator factor as
\[
\begin{aligned}
495535757310206240323260222
 &=2\,3^{16}7^4\,17\,19^{10}\,23,\\
432344158742282808792052775
 &=5^2\,29^7\,139^7.
\end{aligned}
\]
The numerator is not a fifth power.  This proves:

\begin{proposition}\label{prop:quadratic-sector}
The sector $[2,5]$ is empty.
\end{proposition}

\section{The cubic--quartic sector}\label{sec:reducible-end}

\paragraph{Provenance.}
The elimination follows Dahmen--Siksek, Section~7, in particular
Proposition~7.1 and Lemmas~7.2--7.7 \cite[Section~7]{DahmenSiksek}.
Our additional contribution is to certify the ramified local step inside the
intrinsic integral formal neighbourhood, rather than treating model-dependent
coordinates as an unexamined proxy for the formal group.

The quartic direct factor is totally imaginary and unramified outside
$\{3,5,7\}$.  Dahmen--Siksek consider the complete lists of nineteen cubic
fields of signature $(1,1)$ and seventeen quartic fields of signature
$(0,2)$ unramified outside this set.  They compare each cubic--quartic
product with the permitted local algebras at $3$, $5$, and $7$.
Their Proposition~4.3 leaves two products, whose cubic factors are
$\Q(\sqrt[3]{5^\alpha7})$ for $\alpha=1,2$, and whose quartic factor is the
same field in both cases \cite[Proposition~4.3]{DahmenSiksek}.

The surviving field is
\begin{equation}\label{eq:quartic-field}
 M=\Q[a]/(a^4-a^3+3a^2+3a+9)
  =k\!\left(\sqrt{\frac{7+\delta}{2}}\right).
\end{equation}
The associated quotient is the genus-two curve
\begin{equation}\label{eq:section7-curve}
 C_2:\quad V^2=U(U-4)(U^4-4U^3+6U^2+9)
\end{equation}
over $k$.  Write $J_2=\operatorname{Jac}(C_2)$.

\subsection{The Mordell--Weil group and sieve}

We use the Mordell--Weil group determination of
\cite[Lemma~7.3]{DahmenSiksek}, which gives
\begin{equation}\label{eq:section7-MW}
 J_2(k)=\langle D_0\rangle_{10}
       \oplus\Z D_1\oplus\Z D_2\oplus\Z D_3.
\end{equation}
For reproducibility, we also give a direct finite-field certificate for the
remaining $2$-saturation in this statement.  Put
$A=J_2(k)$, $H=\langle D_0,D_1,D_2,D_3\rangle$, and
$H_0=\langle D_0,D_1,D_2,D'_3\rangle$.  The rational and twisted
Mordell--Weil computations used in the cited lemma give
$\rank A=3$; good reduction at $11$ and $13$, together with the displayed
order of $D_0$, gives $A_{\rm tors}=\langle D_0\rangle\simeq\Z/10\Z$.
Moreover, $2A\subseteq H_0\subseteq H$, the second inclusion following from
the verified relation
\[
                 D'_3=9D_0+3D_1-D_2-2D_3.
\]
It therefore remains only to show that the four displayed classes span
$A/2A$.

At the good prime $\mathfrak l=(173,\delta-22)$, the change of variables
\[
 u=-36/U,\qquad v=1296V/U^3
\]
puts the curve into the odd monic form
\[
 v^2=u^5+9u^4+864u^3+28512u^2+373248u+1679616.
\]
Its five roots modulo $\mathfrak l$, in the certified order, are
$15,90,119,122,164$.  Applying the usual odd-degree $u-T$ Kummer map and
then the five quadratic characters of $\F_{173}^{\times}$ to
$D_0,D_1,D_2,D_3$ gives the matrix
\[
 \begin{pmatrix}
 0&0&1&1\\
 1&0&1&1\\
 1&0&1&0\\
 1&1&0&1\\
 1&1&1&1
 \end{pmatrix}.
\]
The determinant of its first four rows is $1$ over $\F_2$.  Hence these
four classes are independent in $A/2A$.  Since
$\dim_{\F_2}A/2A=3+1=4$, they form a basis; as $A/H$ was already killed by
$2$, this proves $A=H$.  The archive contains both the short reconstruction
and an independent evaluation of every Kummer entry.

A Mordell--Weil sieve \cite{BruinStollSieve} with modulus $2240$ yields
twelve points
$P_1,\ldots,P_{12}$ such that every $Q\in C_2(k)$ satisfies
\begin{equation}\label{eq:section7-sieve}
            [Q-P_i]\in2240J_2(k)
\end{equation}
for a unique $i$.

Let $\PP$ be the unique prime of $k$ above $7$, let
$k_\PP$ be the completion, and put $R=\OO_{k_\PP}$, $\m=(\delta)$.  Thus
\[
                         v_\PP(\delta)=1,\qquad v_\PP(7)=2.
\]
Use Flynn's integral $\mathbf P^{15}$ model of $J_2$ with origin
$(1:0:\cdots:0)$ and normalized coordinates $s_i=a_i/a_0$.  On the
origin neighbourhood $\mathcal N$, the projective-coordinate framework of
\cite{CasselsFlynn} and Corollary~2.2 and Theorem~3.5 of \cite{Flynn}
identify $(s_1,s_2)$ as integral formal parameters and give an
integral formal group law.  Define
\begin{equation}\label{eq:G2}
 G[2]=\{A\in\mathcal N:s_1(A),s_2(A)\in\m^2\}.
\end{equation}
It is a subgroup, and the coordinate identities place every nonorigin
normalized coordinate in $\m^2$.

The exact origin-chart valuations of $2240D_1$, $2240D_2$, and $2240D_3$
in the adapted integral coordinates described in Appendix~\ref{app:flynn}
are
\begin{equation}\label{eq:section7-valuations}
\begin{array}{c|l}
D_1&0,2,2,4,4,4,6,6,6,6,8,8,8,8,8,16\\
D_2&0,5,3,10,8,6,15,13,11,9,20,18,16,14,12,20\\
D_3&0,2,2,4,4,4,6,6,6,6,8,8,8,8,8,16.
\end{array}
\end{equation}
The torsion generator is killed by $2240$, so
\begin{equation}\label{eq:2240G2}
                            2240J_2(k)\subset G[2].
\end{equation}

\subsection{The twelve local classes}

For each $P_i$, exact pullback calculations prove
\begin{equation}\label{eq:local-pullback}
 [Q-P_i]\in G[2]\quad\Longrightarrow\quad
                         Q\equiv P_i\pmod{\PP^2}.
\end{equation}
At each of five regular finite bases, the $98$ depth-two congruence classes
split into one base class and $97$ rejected classes.  At each of the five
ramified finite bases, the complete tree is
\[
\begin{array}{c|rrrr|rr}
n&\text{tested}&\text{base}&\text{rejected}&\text{unresolved}
 &49\,\text{-digit lifts}&\text{admissible at }n+1\\ \hline
2&98 &1&91 &6  &294 &294\\
3&294&0&252&42 &2058&686\\
4&686&0&588&98 &4802&686\\
5&686&0&686&0  &--&--.
\end{array}
\]
The last two columns make the branching arithmetic explicit.  Each
unresolved pair modulo $\PP^n$ has $7^2=49$ two-coordinate digit lifts;
the next row counts only those lifts that also satisfy the curve equation
modulo $\PP^{n+1}$.  At the ramified special-fibre point this number is not
uniformly seven, which is why the tested-node column is not obtained by
multiplying the unresolved column by a fixed branching factor.
The singular special-fibre point $U\equiv4\pmod{\PP}$ is included in this
calculation.  At the two points at infinity, the opposite branch lies
outside $\mathcal N$ and the identity branch has
\[
                            s_1=\pm\tfrac12w+O(w^2),
\]
with unit linear coefficient.  Here a node is a full congruence class, and a
rejection is a uniform valuation statement on that class; no finite sample
is used as an exhaustion argument.

Combining \eqref{eq:section7-sieve}, \eqref{eq:2240G2}, and
\eqref{eq:local-pullback}, every $k$-point lies in one of the twelve
explicit $\PP^2$-disks.  Exact $7$-adic and $23$-adic conditions, followed
where necessary by Coleman integration, show that no disk contains a
solution-bearing point beyond its reference $P_i$.  The reference points
themselves either have parameter $0$, $1$, or $\infty$, or fail the required
rational admissibility condition; hence none yields a nonzero primitive
solution.

\begin{proposition}\label{prop:quartic-sector}
The sector $[3,4]$ is empty.
\end{proposition}

\section{The irreducible septic sector}\label{sec:irreducible}

\paragraph{Provenance.}
The seven-field enumeration is Putz's published theorem.  The elimination of
the six pure fields and the exceptional field is the new terminal part of the
argument.

Suppose now that $A_\eta$ is a field.  It has degree $7$, signature
$(1,3)$, is unramified outside $\{3,5,7\}$, and has the prescribed
local algebras at $2$, $3$, $5$, and $7$ of
\cite[Corollary~3.4 and Propositions~3.5--3.7]{DahmenSiksek}.
Putz's exhaustive Hunter enumeration
\cite[Theorem~3.50]{Putz} gives exactly seven possibilities:
\begin{equation}\label{eq:seven-fields}
 \Q\!\left(\sqrt[7]{3^a5}\right)\quad(1\leq a\leq6),
 \qquad
 \Q[T]/(T^7-483T^2+3955T-3945).
\end{equation}
We call the first six fields pure and the last field exceptional.

\subsection{The pure fields}

The Fano-resolvent construction and the rational-parameter theorem in this
subsection supply the geometric obstruction that was absent at the published
seven-field frontier.

The normal closure of a pure field has Frobenius group $F_{42}$ and unique
quadratic subfield
\[
                              K_7=\Q(s),\qquad s^2=-7.
\]

\subsubsection{The Fano-resolvent bridge}

The passage from a pure septic field to a curve is as follows.  Scale
$V=15T$ and write
\begin{equation}\label{eq:scaled-septic}
 h_t(V)=V^7-35V^6+315V^5-15^6t,
 \qquad t=\eta=\frac{X^5}{Z^7}.
\end{equation}
Let $v_1,\ldots,v_7$ be its roots.  A labelled Fano plane $F$ on these
seven roots has seven lines, each a three-element subset; define
\begin{equation}\label{eq:fano-invariant}
 M_F=\sum_{\{i,j,k\}\text{ a line of }F}v_iv_jv_k,
 \qquad q_F=\frac{M_F}{225}.
\end{equation}
There are thirty labelled Fano planes.  They form two $A_7$-orbits
$\mathcal F_+$ and $\mathcal F_-$ of fifteen planes.  Since
\[
 \disc(h_t)=-7^7 15^{36}t^4(t-1)^2,
\]
the two orbit polynomials are conjugate over $K_7$.  After fixing the
Vandermonde orientation, put
\begin{equation}\label{eq:plus-resolvent}
 R_+(q,t)=\prod_{F\in\mathcal F_+}(q-q_F)
          =\frac{U(q,t)+sW(q,t)}2,
 \qquad U,W\in\Z[q,t].
\end{equation}
The normalization of $R_+=0$ is denoted $C_+$.  On its fifteen sheets the
branch permutations above $t=0,1,\infty$ have cycle structures
\[
                         5^3,\qquad3^5,\qquad7^2 1.
\]
Riemann--Hurwitz therefore gives $g(C_+)=3$.

To see why this quotient detects the six pure fields, call
$F\in\mathcal F_+$ and $G\in\mathcal F_-$ compatible when they have no
line in common.  We use unordered pairs $\{F,G\}$, since odd
permutations exchange the two $A_7$-orbits.  Each $F$ has exactly eight
compatible $G$'s, giving 120 such pairs.  The normalization of the
corresponding cover is a curve $X_{42}/\Q$.  The stabilizer of an
unordered compatible pair in $A_7$ is the Frobenius group $F_{21}$,
and its stabilizer in $S_7$ is $F_{42}$.  Thus $X_{42}$ has degree $120$
over the $t$-line and, after base change to $K_7$, the forgetful map
\begin{equation}\label{eq:eight-over-fifteen}
                         X_{42,K_7}\longrightarrow C_+
\end{equation}
has degree $8$.  The branch permutations on the 120 sheets have cycle
structures $5^{24}$, $3^{40}$, and $7^{17}1$, so $g(X_{42})=20$.

If \eqref{eq:scaled-septic} cuts out one of the pure fields
$\Q(\sqrt[7]{3^a5})$, its normal closure has group $F_{42}$.  The
fixed-field/resolvent correspondence supplies a rational point on
$X_{42}$, hence a $K_7$-point on $C_+$ through
\eqref{eq:eight-over-fifteen}.  Both maps preserve $t$.  A nonzero
solution has $t\notin\{0,1,\infty\}$, so the resulting point is not in a
branch fibre.  This constructs the claimed bridge rather than merely
inferring it from matching Galois groups.

The canonical model of $C_+$ is the explicit smooth plane quartic printed
in Appendix~\ref{app:Cplus}.  Its smoothness proves again that its genus is
$3$.  The inverse parameter is the exact rational function
$t=-G_t/H_t$; all coefficients of the two quintics are printed in
Table~\ref{tab:t-parameter-coefficients}.
The archive file \path{p7_f42_canonical_exact.json.gz} specifies
the forward coordinates $N_j(q,t)$ for $j=1,2,3$, and the inverse
forms $G_q,H_q$.  Exact reduction modulo $R_+(q,t)$ verifies
\[
 Q(N_1,N_2,N_3)=0,\qquad
 G_q(N)+qH_q(N)=0,\qquad G_t(N)+tH_t(N)=0,
\]
where $Q$ is the printed quartic and the two inverse denominators are
nonzero in the function field.  Recovering both generators $q$ and $t$
proves that this map induces an isomorphism of function fields, hence of
the smooth projective models, preserving the parameter.

\begin{theorem}[rational-parameter theorem]\label{thm:rational-t}
Let $J_+=\operatorname{Jac}(C_+)$.  Then
\[
 \{Q\in C_+(K_7):t(Q)\in\Pone(\Q)\}
       =\{P_1,P_2,P_3,P_4,P_5\}.
\]
The corresponding parameter values are $0,1,\infty,\infty,\infty$.
\end{theorem}

\begin{proof}
Put $D_i=[P_i-P_1]$.  An exact $2$-descent, implemented in the general
genus-three framework of \cite{BPS}, and a rank calculation give
\[
 \dim_{\F_2}\operatorname{Sel}_2(J_+/K_7)=4,
 \qquad \rank J_+(K_7)=4,
 \qquad \Sha(J_+/K_7)[2]=0,
\]
and $D_2,D_3,D_4,D_5$ form a basis modulo $2$.  Their subgroup therefore
has odd index.  The existence of the fifth-division class used below can
be seen from the branch divisor.  Write $A=(t)_0/5$, a $K_7$-rational
divisor of degree three, and set $B=[A-3P_1]$.  The pole divisor is
$(t)_\infty=7P_3+7P_4+P_5$, so
\[
 5B=7D_3+7D_4+D_5.
\]
Consequently the class $E=3B-4D_3-4D_4$ satisfies
\begin{equation}\label{eq:Erelation}
                         5E=D_3+D_4+3D_5.
\end{equation}
We use $(D_2,E,D_4,D_5)$ for the odd-primary part.  Four
independent local $C_{25}$ rows have determinant $4$ modulo $5$, and
$J_+(K_7)[5]=0$; hence this subgroup is $5$-saturated.

At the split primes
\[
                     23,43,67,71,79,107,109,113,
\]
both reductions of the curve are enumerated, as are the images of every
reduced point in the relevant local quotients of $J_+$.  The same global
coefficient vector must occur at both places.  In addition, rationality of
$t(Q)$ imposes equality of the two projective reductions of $t$.  The
resulting finite sieve has the following exact counts:
\[
\begin{array}{lr}
\toprule
\text{stage}&\text{surviving coefficient classes}\\
\midrule
\text{coupled $2$-primary sieve modulo $16$}&592\\
\text{$5$-primary prefilter modulo $25$}&1252\\
\text{joint point fingerprints modulo $400$}&144\\
\bottomrule
\end{array}
\]
The final classes use exactly the five vectors modulo $25$
\[
 (0,0,0,0),\ (1,0,0,0),\ (0,5,24,22),\
 (0,0,1,0),\ (0,0,0,1),
\]
and their projection at the two primes above $23$ consists precisely of the
five ordered residue pairs of $P_1,\ldots,P_5$.

The combined $23$-adic logarithm matrix has size $6\times4$ and rank $4$.
Two independent annihilating differentials therefore vanish on all of
$J_+(K_7)$.  At the disks of $P_2,P_3,P_4,P_5$, the Siksek tangent
determinants modulo $23$ are
\begin{equation}\label{eq:siksek-dets}
                              12,\quad16,\quad4,\quad19.
\end{equation}
This is the number-field Chabauty criterion of \cite{Siksek}; the values are
units, so each disk contains only its reference point among common
Coleman zeros.

At $P_1$, a normalized second-order system has exactly two simple reduced
zeros,
\[
                         (u,v)=(0,0),\qquad(14,16),
\]
with determinants $20$ and $3$.  The first is $P_1$.  At the second, the two
local parameter values modulo $23^8=78{,}310{,}985{,}281$ are
\[
 t_{s=4}=59{,}233{,}664{,}629,\qquad
 t_{s=19}=24{,}239{,}267{,}738.
\]
Their difference has exact $23$-adic valuation $5$, with first nonzero
quotient digit $9$.  Hence this common Coleman zero does not have rational
$t$.  The five reference points are therefore the entire rational-parameter
locus.
\end{proof}

Each $P_i$ in Theorem~\ref{thm:rational-t} lies above one of the branch
values $0$, $1$, and $\infty$.  A nonzero primitive solution would give a
point with rational parameter outside this set.  Consequently none of the six pure fields occurs.

\subsection{The exceptional field}

The general modularity and conductor theorems used here are published inputs.
The project-specific specialization, exhaustive packet and ray-character
interfaces, and the terminal $q=29$ comparison form the new finite
elimination.

We use the Frey abelian variety of Pacetti and Villagra Torcomian
\cite{PVT} for
\begin{equation}\label{eq:PVT-equation}
                             a^5+b^p+c^3=0.
\end{equation}
In our variables
\begin{equation}\label{eq:PVT-map}
                (a,b,c)=(X,-Z,Y),\qquad
                t_0=-\frac{a^5}{c^3}=\frac{\eta}{\eta-1}.
\end{equation}
Put $F=\Q(\sqrt5)$ and let
\(\bar\rho_7:G_F\to\operatorname{GL}_2(\overline{\F}_7)\) be the residual
representation of their plus motive at $t_0$.  The use of the modular
method at the small exponent $p=7$ rests on the following precise inputs,
not on the headline large-exponent theorem of \cite{PVT}.
Theorem~2.4(1) of \cite{PVT} states that for $r=3$ and $q\geq5$ the plus
motive is modular for every rational specialization.  Its determinant is
the mod-$7$ cyclotomic character (equivalently, the compatible system has
trivial Nebentypus), and Theorem~2.1(3)--(4) gives the needed ramification
control and finiteness at every prime above $7$.  Here $7$ is inert and
unramified in $F$, so there is a single such prime.

Now assume that $\bar\rho_7$ is irreducible.  Apply
\cite[Theorem~7.8]{PVT} with $p=7$ to the specialization
\eqref{eq:PVT-map}.  This theorem packages the modularity, local conductor
calculation, and level-lowering argument for the plus motive: there is a
parallel-weight-two Hilbert newform over $F$, with trivial Nebentypus,
whose residual representation is $\bar\rho_7$ and whose level is one of
\begin{equation}\label{eq:four-levels}
 3^{\epsilon_3}(\sqrt5)^{\epsilon_5},\qquad
 (\epsilon_3,\epsilon_5)\in\{2,3\}^2.
\end{equation}

The conclusion here is a newform conclusion, but the finite elimination below
does not depend on reconstructing a new quotient.  Let $V_N$ be the full
mod-$7$ cuspidal Brandt module at one of the four levels~$N$ in
\eqref{eq:four-levels}.  The historical computation first imposed eighteen
local trace-polynomial conditions, so its output must not be called the full
ambient module.  We instead verify the coverage of every condition.

Section~7.4 of \cite{PVT} gives an exhaustive trichotomy at an auxiliary
prime: ordinary specialization, one of the two degenerations $t=0,\infty$,
or the level-lowering specialization $t=1$.  From the complete upstream
\texttt{Data.txt}, the certificate independently reduces all $6253$ trace
polynomials at $37$ primes modulo~$7$, verifies all ordinary parameters and
all five and three degenerate character values, performs the required
relative-degree-two trace transport, and adjoins both signs of
$N(\mathfrak q)+1$.  Seventeen of the eighteen historical filters are exactly
the resulting squarefree allowed-trace polynomials.  At $\ell=131$ the
historical degree-$48$ filter omitted the allowed level-lowering trace $1$,
since $-(131+1)\equiv1\pmod7$.  We therefore replace only that condition, in
memory, by $T^{49}-T$, which admits every $\F_{49}$-valued trace.

Write $W_N\subseteq V_N$ for the intersection after these seventeen verified
conditions and the conservative no-filter condition at $131$.  Characteristic-zero Jacquet--Langlands
\cite[Theorem~3.9 and \S4]{DembeleVoight} supplies the corresponding
Hecke eigenvector in the definite quaternionic function model.  The exact
order and neighbour computations verify the maximal order, class number
one, and the complete Hecke operators; all finite stabilizer orders divide
$120$, hence are prime to $7$.  The integral function lattice and its
cuspidal mass condition are therefore preserved at primes above $7$.
Scale the eigenvector to be primitive in this local lattice and reduce.
The resulting nonzero vector has the required residual eigenvalues.
The local trichotomy places it in $W_N$.  This proves coverage without
degeneracy-map subtraction.  A fresh simultaneous primary decomposition has
packet dimensions, with multiplicity, summing to $\dim W_N$ at every level.
The prime-$2$ calculation gives the necessary residual trace $-1$ when $X$ is
even and $0$ when $Y$ is even.  Retaining every packet for which that scalar
is a root of the residual $T_2$ polynomial gives
\[
\begin{array}{c|rrrr}
 &22&23&32&33\\ \hline
 \dim W_N&8&38&28&98\\
 \text{all packets}&4&14&12&24\\
 X\text{-even}&1&2&2&5\\
 Y\text{-even}&2&10&3&11.
\end{array}
\]
Across all thirty-six retained packet occurrences, the set of residual
$T_{\mathfrak q}$ annihilators at $\mathfrak q\mid29$ is exactly
\[
 T,\quad T^2,\quad T+2,\quad T^2+5T+2,\quad T^2+3T+4.
\]
The terminal list already contains all five, because $T+2=T-5$ modulo~$7$;
together with $T-2$ it is the same six-factor list used in the resultant
calculation below.  Thus every residual representation supplied by level
lowering survives the preceding filters and is covered before the terminal
comparison.  The old/new calculation remains only a diagnostic.  We make no
complete-newspace claim and require no mod-$7$ injectivity or semisimplicity
statement for degeneracy maps.

Suppose instead that $\bar\rho_7$ is reducible.  Its semisimplification has
the determinant-compatible form
\begin{equation}\label{eq:reducible-semisimplification}
 \bar\rho_7^{\mathrm{ss}}
       =\psi\oplus\psi^{-1}\bar\chi_7.
\end{equation}
At a finite prime $\lambda\nmid7$, the cyclotomic character is unramified,
so the two characters in \eqref{eq:reducible-semisimplification} have the
same Artin-conductor exponent.  The conductor of the semisimplification
is at most that of the representation.  Thus the local conductor bounds
at the primes above $3$ and $5$ give
\[
 2a_\lambda(\psi)
 =a_\lambda(\bar\rho_7^{\mathrm{ss}})
 \leq a_\lambda(\bar\rho_7)\leq3,
 \qquad a_\lambda(\psi)\leq1.
\]
At every other finite prime away from $7$, the character $\psi$ is
unramified.

The prime $\mathfrak p=(7)$ is unramified and inert in $F$, with residue
field $\F_{49}$.  By the finite-flat character classification
\cite[Corollary~3.4.4]{Raynaud1974}, its restriction to inertia has the form
\[
 \psi|_{I_{\mathfrak p}}=\omega_2^{r_0+7r_1},
 \qquad r_0,r_1\in\{0,1\},
\]
The character has level dividing two because it is a rank-one character
of the full local Galois group.  Here $\omega_2$ is a fundamental
character of level two, normalized so
that $\omega_2^{8}=\bar\chi_7|_{I_{\mathfrak p}}$.
The two mixed possibilities must be excluded globally.  Put
$\varepsilon=(1+\sqrt5)/2$ and
\[
 u=\varepsilon^8=13+21\varepsilon.
\]
The unit $u$ is totally positive and satisfies
\[
 u\equiv1\pmod{3(\sqrt5)},\qquad u\equiv-1\pmod{7}.
\]
Global reciprocity applied to $u$ shows that the value of $\psi$ on its
local reciprocity image at $\mathfrak p$ is $1$: at every other finite
place this follows from the conductor bounds just proved, and at both
real places it follows from total positivity.  The displayed inertial
character gives that value as $(-1)^{r_0+r_1}$.  Hence $r_0=r_1$.
If both are zero then $\psi$ is unramified at $\mathfrak p$; if both are
one then $\psi^{-1}\bar\chi_7$ is unramified there.  Interchanging the two
characters if necessary, we may therefore choose $\psi$ such that
\begin{equation}\label{eq:character-conductor}
 \operatorname{cond}(\psi)\mid3(\sqrt5).
\end{equation}
Including both real places gives a modulus containing the conductor of
$\psi$ regardless of its signs.  Its narrow ray group is
$C_4\times C_2$, so there are eight possible characters, all of which are
included in the comparison.

At $q=29$, take
$\mathfrak q=(29,\sqrt5-11)$; the conjugate prime gives the conjugate
packet and the same common-root test.  The exceptional septic polynomial
is squarefree modulo $29$ with factor degrees $[1,3,3]$.  Among
$\eta\in\F_{29}\setminus\{0,1\}$, the polynomial $\Phi(T)-\eta$ has
this pattern precisely for $\eta=10,14,24,28$.

The branch residues do not give this pattern.  If $29\mid X$, scale the
five-root cluster at $T=0$ using $\ord_{29}(\eta)=5\ord_{29}(X)$.
Its separable reduction has degrees $[1,2,2]$, since $29\equiv-1\pmod5$;
the complementary quadratic has degrees $[1,1]$.  If $29\mid Y$, the
corresponding three-root cluster at $T=1$ has degrees $[1,2]$, and the
complementary quartic has degrees $[1,1,2]$.  Both cases therefore have
pattern $[1,1,1,2,2]$.  If $29\mid Z$, the scaled coordinate $U=ZT$
gives reduction $15U^7-X^5$; its pattern is $[1,1,1,1,1,1,1]$ or $[7]$
because $\F_{29}$ contains the seventh roots of unity.  Hensel's lemma
applies to each scaled separable reduction.  Thus only the four ordinary
values remain.  By \eqref{eq:PVT-map} they correspond to
$t_0=14,10,25,15$.  Direct point counting on
\[
                y^2=5x^6-12x^5+10t_0x^3+t_0^2
\]
gives the following polynomials over $\F_7$ satisfied by the two
conjugate real-multiplication Frobenius traces.  Explicitly, if
$N_1=\#C_{t_0}(\F_{29})$, $N_2=\#C_{t_0}(\F_{29^2})$,
$A=30-N_1$, and $B=(A^2-(842-N_2))/2$, the degree-four Frobenius
polynomial is $X^4-AX^3+BX^2-29AX+29^2$, and the two traces satisfy
$T^2-AT+(B-58)$.  The table records the latter polynomial modulo $7$:
\begin{equation}\label{eq:curve-polynomials}
\begin{array}{c|c|c}
\eta&t_0&\text{curve polynomial}\\ \hline
10&14&T^2+6T+6\\
14&10&T^2+6T+4\\
24&25&T^2+4\\
28&15&T^2+2T+3.
\end{array}
\end{equation}
For clarity, the next list is an \emph{annihilating family}; its members are
not asserted to be irreducible or squarefree.  For each of the four ambient
cuspidal spaces, reduce the characteristic polynomial of the Hecke operator
$T_{\mathfrak q}$ modulo $7$ on every packet surviving the conservative
$T_2$ condition.  The preceding dimension accounting proves that every
residual eigenvalue $a_{\mathfrak q}(f)\bmod\mathfrak p$, for every relevant
newform and each coefficient prime $\mathfrak p\mid7$ realizing the required
Frey congruence, is a root of
one of the first four or the sixth displayed factors.  We retain the factors as they
occur in the packet calculation, so they need not be squarefree or minimal.
The reducible characters contribute the possible traces
$\psi(\mathfrak q)+\psi(\mathfrak q)^{-1}$.  The resulting complete
annihilating family is
\begin{equation}\label{eq:test-polynomials}
 T,\quad T^2,\quad T^2+5T+2,\quad T^2+3T+4,\quad T-2,\quad T-5.
\end{equation}
Here $T^2$ and $T^2+3T+4=(T-2)^2$ deliberately record multiplicity; for
the resultant test they have the same roots as $T$ and $T-2$.  For the specified prime $\mathfrak q$, the ray class has order two:
a totally positive generator is $6-\varepsilon$, whose square is
congruent to the totally positive unit $\varepsilon^2$ modulo
$3(\sqrt5)$; its class is nontrivial, since modulo $3$ the generator is
$-\varepsilon=\varepsilon^5$, whereas every totally positive unit is an
even power of $\varepsilon$, whose order modulo $3$ is $8$.
Thus the reducible traces are
$2$ and $-2$, covered by the displayed linear factors.  The
$24$ pairwise resultants of \eqref{eq:curve-polynomials} and
\eqref{eq:test-polynomials} form the matrix
\begin{equation}\label{eq:resultants}
 \begin{pmatrix}
 6&1&5&1&1&5\\
 4&2&3&1&6&3\\
 4&2&6&1&1&1\\
 3&2&6&2&4&3
 \end{pmatrix}\pmod7.
\end{equation}
Every entry is nonzero.  Thus the exceptional field is incompatible with
both the irreducible and reducible residual alternatives.

\begin{proposition}\label{prop:septic-sector}
The sector $[7]$ is empty.
\end{proposition}

\section{Proof of the main theorem}

\begin{proof}[Proof of Theorem~\ref{thm:main}]
Suppose that a nonzero primitive solution exists and reorder it as in
\eqref{eq:DS-order}.  The descent algebra $A_\eta$ belongs to exactly one
of the four sectors in \eqref{eq:four-sectors}.  Proposition
\ref{prop:rational-sector} excludes $[1,6]$.  Proposition
\ref{prop:quadratic-sector} excludes $[2,5]$.  Proposition
\ref{prop:quartic-sector} excludes $[3,4]$.  Proposition
\ref{prop:septic-sector} excludes $[7]$.  Every possible factorization is
therefore impossible, a contradiction.
\end{proof}

\section{Computational certification}

The computer-assisted portions of the proof are finite statements over
$\Q$, number fields, finite fields, and truncated local rings.  They are
organized so that discovery is not used as evidence: every accepted output
is checked against exact input data and against the mathematical predicate
stated in the relevant proposition.  The certificate archive serves two
different purposes.  It independently reconstructs finite calculations from
the published part of the proof, and it supplies the load-bearing arithmetic
for the new seven-field elimination.  The provenance column below keeps
those purposes distinct.

\subsection{Exact arithmetic layer}

The exact layer checks polynomial identities, resultants, discriminants,
field embeddings, ideal-class representatives, divisor relations, point
incidence, and rational substitutions.  Integer and rational calculations
are repeated without floating-point arithmetic.  Manifests identify the
input bytes and the corresponding outputs.

\subsection{Finite exhaustion layer}

Whenever a finite search is used as an exhaustion, the searched set is
defined before enumeration.  This applies to local factor patterns, Selmer
classes, reduced curve points, Mordell--Weil coefficient classes, residue
trees, ray characters, newform packets, and resultant comparisons.  A
residue-tree node denotes a complete congruence class.  A rejected node is
accompanied by a uniform congruence or valuation obstruction; it is not a
sample point.

\subsection{Local analytic layer}

Formal-group and Coleman calculations record their working precision,
unit pivots, ranks, determinants, and higher-precision replays.  Divisions
are made only by certified units or with their exact loss of precision
recorded.  Simple-root conclusions are accompanied by a nonzero reduced
Jacobian determinant.  Singular classes are subdivided or normalized until
the stated conclusion follows.

\subsection{Software and trust boundary}

The calculations use PARI/GP \cite{PARI} for exact number-field arithmetic,
SageMath \cite{Sage} for curve and finite-field calculations, and Magma
\cite{Magma} for descent,
Mordell--Weil, saturation, newform, and Chabauty routines.  Their role is
the same as that of large integer arithmetic or a table of modular forms in
a conventional proof: the inputs, outputs, and verification conditions are
part of the mathematical record.  The argument also imports the general
theorems cited in the bibliography, rather than attempting to reprove them
from first principles.

\subsection{Companion certificate archive}

The companion archive preserves the evidence in repository-relative paths
and provides a single integrity and interface verifier.  Running
\begin{center}
\texttt{python3 -B scripts/verify\_dependency\_closed\_release\_v3.py}
\end{center}
from the archive root checks the declared release manifest, the $52$ base
evidence manifests, the exact $765$-file repository closure, the explicit
added-file and generator-input indexes, and consistency of the
theorem-composition graph.  The optional pair of flags
\texttt{--run-sector-replays --run-foundational-replays} reruns both the ten
sector entry points and the foundational reconstruction: the two independent
five-place workers, their semantic comparison, the Selmer projector and
rank-four deduction, both split-$23$ global logarithm matrices, the combined
annihilator, and the exceptional local calculation.  The distributed complete
log was produced from a clean extraction in an isolated filesystem in which
the original workspace path was absent.  In the arXiv source bundle the
release is the ancillary file
\path{anc/PRIMITIVE_357_REPRODUCIBILITY_V3.tar.gz}; the accompanying
\texttt{PROGRAMS\_AND\_CERTIFICATES.md} gives proposition-by-proposition
commands and paths.  The principal correspondence is:
\begin{center}
\begin{tabular}{@{}>{\raggedright\arraybackslash}p{.22\textwidth}
                    >{\raggedright\arraybackslash}p{.20\textwidth}
                    >{\raggedright\arraybackslash}p{.48\textwidth}@{}}
\toprule
mathematical assertion & provenance & certificate family\\
\midrule
four exhaustive factor sectors
  & published; reconstructed
  & signed global-algebra superselection\\
Proposition~\ref{prop:rational-sector}
  & published; reconstructed
  & rational-factor composition and its Mordell--Weil, sieve, and local packages\\
Proposition~\ref{prop:quadratic-sector}
  & published, with new intrinsic routing
  & database-free quadratic router and quadratic-factor reconstruction\\
Proposition~\ref{prop:quartic-sector}
  & published, with new formal-group certification
  & quartic router, Mordell--Weil basis, sieve, and intrinsic formal-group packages\\
Theorem~\ref{thm:rational-t} and Appendix~\ref{app:Cplus}
  & new
  & Fano-cover, canonical-quartic, rational-parameter, and split-prime packages\\
exceptional-field elimination
  & new, using published modular theorems
  & Putz--PVT interface and finite $q=29$ comparison packages\\
Proposition~\ref{prop:septic-sector}
  & new terminal composition
  & irreducible-sector composition package\\
Theorem~\ref{thm:main}
  & new global composition
  & complete four-sector proof graph and terminal certificate\\
\bottomrule
\end{tabular}
\end{center}
The archive distinguishes three operations.  The default verifier is a
dependency-free check of identity, coverage, and exact finite interfaces.
The supplied Python and PARI/GP programs rerun the open exact-arithmetic
layers.  SageMath and Magma programs rerun the curve, descent, modular-form,
and Chabauty layers when those systems are available; authenticated logs are
included for computations requiring the proprietary Magma runtime.  Imported
published theorems remain bibliographic dependencies and are not represented
as newly proved computational claims.

\section{Concluding perspective}

The proof is not driven by the size of a putative solution.  It represents
that solution simultaneously by a descent algebra, local factor labels,
divisor classes, finite quotient coordinates, and $p$-adic analytic data.
Each representation imposes an exact compatibility condition.  The global
contradiction comes from the empty intersection of those conditions.

This organization is useful beyond the present signature.  It separates
three tasks that are often entangled: enumerating the possible algebraic
states, transporting a state to a curve where a finitely generated group
acts, and testing the remaining local states with analytic functions.  It
also makes the proof modular: a general theorem can be replaced by an
independent implementation without changing the surrounding argument, while
the finite interfaces remain explicit.

\appendix

\section{The plane quartic in the pure septic case}\label{app:Cplus}

Let $K_7=\Q(s)$ with $s^2=-7$.  In the monomial order
\[
 x^4,x^3y,x^3z,x^2y^2,x^2yz,x^2z^2,xy^3,xy^2z,xyz^2,xz^3,
 y^4,y^3z,y^2z^2,yz^3,z^4,
\]
write the coefficient of each monomial as $(U+sW)/2$.  The pairs $(U,W)$
for $C_+$ are recorded below with the monomial repeated on every row, so
that each coefficient can be quoted and checked independently of the
surrounding order.
\begin{landscape}
\small
\renewcommand{\arraystretch}{1.22}
\begin{longtable}{@{}c>{\raggedleft\arraybackslash}p{82mm}>{\raggedleft\arraybackslash}p{82mm}@{}}
\caption{Exact coefficients of the plane quartic $C_+$.  The coefficient
of $m$ is $(U+sW)/2$.}\label{tab:cplus-quartic-coefficients}\\
\toprule
monomial $m$ & $U$ & $W$\\
\midrule
\endfirsthead
\multicolumn{3}{c}{\tablename\ \thetable\ (continued)}\\
\toprule
monomial $m$ & $U$ & $W$\\
\midrule
\endhead
$x^4$       & $526589143431946875/69798278$   & $-114868376351859375/69798278$\\
$x^3y$      & $-1271586888400500/3172649$     & $-870445613016000/3172649$\\
$x^3z$      & $72385949750625/34899139$        & $-402884946883125/34899139$\\
$x^2y^2$    & $-666163656476955/34899139$      & $264021753009735/34899139$\\
$x^2yz$     & $-74662865230275/34899139$       & $45051674372775/34899139$\\
$x^2z^2$    & $-20453292214425/69798278$       & $1113321269025/69798278$\\
$xy^3$      & $77961057423858/174495695$       & $1554848756622/34899139$\\
$xy^2z$     & $5974117172919/34899139$         & $534487214709/34899139$\\
$xyz^2$     & $251941629933/34899139$          & $209684218779/34899139$\\
$xz^3$      & $-13800316635/69798278$          & $9656077725/69798278$\\
$y^4$       & $-11467163171457/8724784750$     & $-8293918213539/8724784750$\\
$y^3z$      & $-192270890133/174495695$        & $-111061125567/174495695$\\
$y^2z^2$    & $-4669897617/31726490$           & $-23997395583/158632450$\\
$yz^3$      & $4101243048/174495695$           & $-629073270/34899139$\\
$z^4$       & $2$                              & $0$\\
\bottomrule
\end{longtable}
\end{landscape}
The equation and its three first partial derivatives have no common
projective zero over $K_7$; this is checked by exact elimination.  As an
independent good-reduction certificate, both reductions at the primes above
$23$ are smooth plane quartics.  Thus $C_+$ is smooth, nonhyperelliptic, and
has genus $3$.  The two reductions and their common zeta numerator are
recorded in Appendix~\ref{app:p23-data}.
The five points in Theorem~\ref{thm:rational-t} are
\begin{align*}
P_1={}&[-113-139s:32925+11475s:-138600-57960s],\\
P_2={}&[1109-305s:4325-19625s:163800-37800s],\\
P_3={}&[-1+s:50:0],\\
P_4={}&[5+s:400:0],\\
P_5={}&[-157167-154941s:-82225175+2993875s:186832800].
\end{align*}
The homogeneous forms $G_t$ and $H_t$ have degree $5$.  For each monomial
$m=x^ay^bz^c$ of degree $5$, write
\[
 [m]G_t=\frac{U_{G,m}+sW_{G,m}}2,
 \qquad
 [m]H_t=\frac{U_{H,m}+sW_{H,m}}2.
\]
Table~\ref{tab:t-parameter-coefficients} gives the four rational components
for every monomial; a component is displayed as separate numerator and
positive denominator.  These data define the parameter completely, with no
external coefficient file required.

% BEGIN INLINED EXACT CPLUS T PARAMETER TABLE
% These entries are evidence, not decorative data.  Each exact fraction needs
% enough vertical room for its numerator and denominator to remain visually
% distinct at ordinary reading magnification.
% Fixed-size ledger leaves avoid the fragile automatic splitting of unusually
% tall fraction rows.  Every continuation page therefore has its own heading
% and contains complete coefficient pairs.
\newcommand{\ExactLedgerBare}{%
  \begin{center}%
  \begin{tabular}{@{}cc>{\raggedright\arraybackslash}p{170mm}@{}}}
\newcommand{\ExactLedgerHead}{%
  \ExactLedgerBare
  \toprule
  monomial & component & rational value\\
  \midrule}
\newcommand{\ExactLedgerTail}{%
  \bottomrule
  \end{tabular}%
  \end{center}}
\newcommand{\ExactLedgerNext}{%
  \ExactLedgerTail
  \end{landscape}
  \begin{landscape}
  \small
  \renewcommand{\arraystretch}{2.5}
  \begin{center}\tablename\ \ref{tab:t-parameter-coefficients} (continued)\end{center}
  \ExactLedgerHead}
\begin{landscape}
\small
\renewcommand{\arraystretch}{2.5}
\begin{center}
\captionof{table}{Exact coefficient ledger for the inverse parameter $t=-G_t/H_t$.
For each monomial and each form, the two successive rows are the rational
components $U$ and $W$ of the coefficient $(U+sW)/2$.}
\label{tab:t-parameter-coefficients}
\end{center}
\ExactLedgerHead
\multicolumn{3}{c}{$G_t$}\\
$x^{5}$ & $U$ & $\frac{\mathtt{8724124647223604676837903272778053868067651476251979316420946130788203125}}{\mathtt{1262470913738019638514233667224982865889496755254608725372117056}}$\\
$x^{5}$ & $W$ & $\frac{\mathtt{104539127651167761178849503246013584932810386714540106341619861276432421875}}{\mathtt{34086714670926530239884309015074537379016412391874435585047160512}}$\\
$x^{4}y$ & $U$ & $\frac{\mathtt{160516210802709714993563463849330968569737125845236059453160939760278125}}{\mathtt{516465373801917124846731954773856626954794127149612660379502432}}$\\
$x^{4}y$ & $W$ & $\frac{\mathtt{-1516029210727323321897909052657899390133429038437597834516845689254996875}}{\mathtt{4648188364217254123620587592964709642593147144346513943415521888}}$\\
$x^{4}z$ & $U$ & $\frac{\mathtt{-3431261016436648118361864834749807161605723145608327043665883055988296875}}{\mathtt{536865756067092851278177866987423963719508495172022360464492778064}}$\\
$x^{4}z$ & $W$ & $\frac{\mathtt{-14405986606016247091870808274044492866667442818955446375070000269331640625}}{\mathtt{536865756067092851278177866987423963719508495172022360464492778064}}$\\
$x^{3}y^{2}$ & $U$ & $\frac{\mathtt{-156765051716182227120181782804196796183006994771851219874584685103262375}}{\mathtt{6391259000798724419978307940326475758565577323476456672196342596}}$\\
$x^{3}y^{2}$ & $W$ & $\frac{\mathtt{-2797425838906852452662912431708042341593603601309085420184135364665875}}{\mathtt{3195629500399362209989153970163237879282788661738228336098171298}}$\\
\ExactLedgerNext
$x^{3}yz$ & $U$ & $\frac{\mathtt{-2971418364998561881697119480380162120182637893006270782575715726325204375}}{\mathtt{805298634100639276917266800481135945579262742758033540696739167096}}$\\
$x^{3}yz$ & $W$ & $\frac{\mathtt{581935828601695609265701070325316097436069202977166807199027638486990625}}{\mathtt{805298634100639276917266800481135945579262742758033540696739167096}}$\\
$x^{3}z^{2}$ & $U$ & $\frac{\mathtt{-1101411663992169427898989261604416930376364487448681819065769508346053125}}{\mathtt{11274180877408949876841735206735903238109678398612469569754348339344}}$\\
$x^{3}z^{2}$ & $W$ & $\frac{\mathtt{-107499314624920129138097362917410399320732876511214060737904168396284375}}{\mathtt{3758060292469649958947245068911967746036559466204156523251449446448}}$\\
$x^{2}y^{3}$ & $U$ & $\frac{\mathtt{456938855157174881485528038353970608056813975973606988967508412987185}}{\mathtt{4260839333865816279985538626884317172377051548984304448130895064}}$\\
$x^{2}y^{3}$ & $W$ & $\frac{\mathtt{1661711353625849642579873501624741967233922652605715624836234484805605}}{\mathtt{12782518001597448839956615880652951517131154646952913344392685192}}$\\
$x^{2}y^{2}z$ & $U$ & $\frac{\mathtt{8474129667022123046619324483315222687658910122135070504455178101736425}}{\mathtt{115042662014377039559609542925876563654180391822576220099534166728}}$\\
$x^{2}y^{2}z$ & $W$ & $\frac{\mathtt{10719261028668655910991493734095990759234174142059301575470250588670875}}{\mathtt{268432878033546425639088933493711981859754247586011180232246389032}}$\\
\ExactLedgerNext
$x^{2}yz^{2}$ & $U$ & $\frac{\mathtt{4788522504016908727437175426249647578416822156275084241334728620014875}}{\mathtt{16911271316113424815262602810103854857164517597918704354631522509016}}$\\
$x^{2}yz^{2}$ & $W$ & $\frac{\mathtt{5591416257739747140557321332572949584770401809552651792943380516721625}}{\mathtt{1879030146234824979473622534455983873018279733102078261625724723224}}$\\
$x^{2}z^{3}$ & $U$ & $\frac{\mathtt{-673613683718293366731089090657452594859133231945568100511338992953086875}}{\mathtt{4261640371660583053446175908146171424005458434675513497367143672272032}}$\\
$x^{2}z^{3}$ & $W$ & $\frac{\mathtt{727154975549102369129278091212463518913101596132005437190890995061885625}}{\mathtt{4261640371660583053446175908146171424005458434675513497367143672272032}}$\\
$xy^{4}$ & $U$ & $\frac{\mathtt{181934888585282020021801499997708888564901966896834717514388222336671}}{\mathtt{127825180015974488399566158806529515171311546469529133443926851920}}$\\
$xy^{4}$ & $W$ & $\frac{\mathtt{-66576524945702734300367034318330080309880903641645008587735355126139}}{\mathtt{127825180015974488399566158806529515171311546469529133443926851920}}$\\
$xy^{3}z$ & $U$ & $\frac{\mathtt{55752557283144284245053396819690398445556264921670780430266851605377}}{\mathtt{134216439016773212819544466746855990929877123793005590116123194516}}$\\
$xy^{3}z$ & $W$ & $\frac{\mathtt{-154753525348486993811829134912142488930864130682803774941942077141595}}{\mathtt{402649317050319638458633400240567972789631371379016770348369583548}}$\\
\ExactLedgerNext
$xy^{2}z^{2}$ & $U$ & $\frac{\mathtt{45466049893188675406917805349188754066362340943964804016941934609315}}{\mathtt{1537388301464856801387509346373077714287683417992609486784683864456}}$\\
$xy^{2}z^{2}$ & $W$ & $\frac{\mathtt{-90644573939421031699724276084922394905211688255008129426169575178615}}{\mathtt{1537388301464856801387509346373077714287683417992609486784683864456}}$\\
$xyz^{3}$ & $U$ & $\frac{\mathtt{86176315804162437100883836603358200621380221474103257474111818725307525}}{\mathtt{6392460557490874580169263862219257136008187652013270246050715508408048}}$\\
$xyz^{3}$ & $W$ & $\frac{\mathtt{-14901625678223850735415766370606467751263679923368054621027598776871975}}{\mathtt{6392460557490874580169263862219257136008187652013270246050715508408048}}$\\
$xz^{4}$ & $U$ & $\mathtt{0}$\\
$xz^{4}$ & $W$ & $\mathtt{0}$\\
$y^{5}$ & $U$ & $\frac{\mathtt{-50208669906374986341132909015146131126007030919811841716909}}{\mathtt{18313514842869689191983526987088351258653049645369350436400}}$\\
$y^{5}$ & $W$ & $\frac{\mathtt{-118965057882822154213596381597899423461490949136300188722323}}{\mathtt{91567574214348445959917634935441756293265248226846752182000}}$\\
\ExactLedgerNext
$y^{4}z$ & $U$ & $\frac{\mathtt{-317869767610131557388806079392661467805959663406188675965981}}{\mathtt{164821633585827202727851742883795161327877446808324153927600}}$\\
$y^{4}z$ & $W$ & $\frac{\mathtt{-20431957573101268107747459729242315992574863384255199500651}}{\mathtt{32964326717165440545570348576759032265575489361664830785520}}$\\
$y^{3}z^{2}$ & $U$ & $\frac{\mathtt{-90218411306089523014211143406416603450262335284699082090763}}{\mathtt{346125430530237125728488660055969838788542638297480723247960}}$\\
$y^{3}z^{2}$ & $W$ & $\frac{\mathtt{-6384772760573948911903427720751451666967418352649934952949}}{\mathtt{69225086106047425145697732011193967757708527659496144649592}}$\\
$y^{2}z^{3}$ & $U$ & $\frac{\mathtt{34563607476891939880909708193414566015508357013210019043845}}{\mathtt{13083541274042963352536871350115659906206911727644771338772888}}$\\
$y^{2}z^{3}$ & $W$ & $\frac{\mathtt{-150431184828356847118153754501573195221720737251858300004385}}{\mathtt{13083541274042963352536871350115659906206911727644771338772888}}$\\
$yz^{4}$ & $U$ & $\mathtt{0}$\\
$yz^{4}$ & $W$ & $\mathtt{0}$\\
$z^{5}$ & $U$ & $\mathtt{0}$\\
$z^{5}$ & $W$ & $\mathtt{0}$\\
\ExactLedgerTail
\end{landscape}
\begin{landscape}
\small
\renewcommand{\arraystretch}{2.5}
\ExactLedgerBare
\multicolumn{3}{c}{\tablename\ \ref{tab:t-parameter-coefficients} (continued)}\\
\toprule
monomial & component & rational value\\
\midrule
\multicolumn{3}{c}{$H_t$}\\
$x^{5}$ & $U$ & $\frac{\mathtt{417540987588276327046139642388996781228921135238823876465270116076953125}}{\mathtt{55097060351713141551537137416607549642806701064452212691347781}}$\\
$x^{5}$ & $W$ & $\frac{\mathtt{271400501016398695076663568050301774206517949977640898218621509239453125}}{\mathtt{55097060351713141551537137416607549642806701064452212691347781}}$\\
$x^{4}y$ & $U$ & $\frac{\mathtt{40800518279237110617651643358322019052974596058651349157620400597734375}}{\mathtt{40070589346700466582936099939350945194768509865056154684616568}}$\\
$x^{4}y$ & $W$ & $\frac{\mathtt{-10475518577744892148533246091047638248959642317009125512752949147203125}}{\mathtt{40070589346700466582936099939350945194768509865056154684616568}}$\\
$x^{4}z$ & $U$ & $\frac{\mathtt{295466106564310819299307055245131237369729837129047525997756150038046875}}{\mathtt{13884459208631711670987358628985102509987288668241957598219640812}}$\\
$x^{4}z$ & $W$ & $\frac{\mathtt{-97401726931434571033238656526295780061095710676931455291585333023515625}}{\mathtt{13884459208631711670987358628985102509987288668241957598219640812}}$\\
$x^{3}y^{2}$ & $U$ & $\frac{\mathtt{-2408541684482196326060738206798781548127364231879824910739247944739375}}{\mathtt{73462747135617522068716183222143399523742268085936283588463708}}$\\
$x^{3}y^{2}$ & $W$ & $\frac{\mathtt{-1401617439695752490670177104164771958374512006836961227347342807905625}}{\mathtt{220388241406852566206148549666430198571226804257808850765391124}}$\\
\ExactLedgerNext
$x^{3}yz$ & $U$ & $\frac{\mathtt{-88727646196045391066194147279337983419623095543345635287685584174321875}}{\mathtt{13884459208631711670987358628985102509987288668241957598219640812}}$\\
$x^{3}yz$ & $W$ & $\frac{\mathtt{-294670882430653452222330070241389024327059653868933829299253113096875}}{\mathtt{661164724220557698618445648999290595713680412773426552296173372}}$\\
$x^{3}z^{2}$ & $U$ & $\frac{\mathtt{-143378643984727743099590687262597995420899449430994725195546663512734375}}{\mathtt{1166294573525063780362938124834748610838932248132324438250449828208}}$\\
$x^{3}z^{2}$ & $W$ & $\frac{\mathtt{-18365708170095511831980056438251994847190327220498321285191544780928125}}{\mathtt{166613510503580540051848303547821230119847464018903491178635689744}}$\\
$x^{2}y^{3}$ & $U$ & $\frac{\mathtt{8623734371967307823196434999403662938951554222083721534345295726825}}{\mathtt{110194120703426283103074274833215099285613402128904425382695562}}$\\
$x^{2}y^{3}$ & $W$ & $\frac{\mathtt{18881132466481791978253130146174121549567643122623579546749820507925}}{\mathtt{110194120703426283103074274833215099285613402128904425382695562}}$\\
$x^{2}y^{2}z$ & $U$ & $\frac{\mathtt{28562333150244501488893362369966137889574552762163770883392632257375}}{\mathtt{661164724220557698618445648999290595713680412773426552296173372}}$\\
$x^{2}y^{2}z$ & $W$ & $\frac{\mathtt{127302955155020767473662087940112338634670901473338060096039562505625}}{\mathtt{1983494172661673095855336946997871787141041238320279656888520116}}$\\
\ExactLedgerNext
$x^{2}yz^{2}$ & $U$ & $\frac{\mathtt{111047168050090146883232536414177868796628159548247958136634257244625}}{\mathtt{583147286762531890181469062417374305419466124066162219125224914104}}$\\
$x^{2}yz^{2}$ & $W$ & $\frac{\mathtt{808931698949954129585743924035840186651521792034283998632202423247375}}{\mathtt{194382428920843963393823020805791435139822041355387406375074971368}}$\\
$x^{2}z^{3}$ & $U$ & $\frac{\mathtt{56671479059928976818861058588328919604601251570358478252116683982508125}}{\mathtt{440859348792474108977190611187534974897116389794018637658670035062624}}$\\
$x^{2}z^{3}$ & $W$ & $\frac{\mathtt{71221107904627237798651209022651158397866104168849046635811093612975625}}{\mathtt{440859348792474108977190611187534974897116389794018637658670035062624}}$\\
$xy^{4}$ & $U$ & $\frac{\mathtt{15345599112192921388150576223890339236281253816603221545634911025}}{\mathtt{18365686783904380517179045805535849880935567021484070897115927}}$\\
$xy^{4}$ & $W$ & $\frac{\mathtt{-32203376966478498296624374226292138386494075379549015420791430453}}{\mathtt{55097060351713141551537137416607549642806701064452212691347781}}$\\
$xy^{3}z$ & $U$ & $\frac{\mathtt{3144967508609794715796497440988774261265066246157112647816548077300}}{\mathtt{3471114802157927917746839657246275627496822167060489399554910203}}$\\
$xy^{3}z$ & $W$ & $\frac{\mathtt{-1516919058667068918402436722904874293272390262093914022875222741380}}{\mathtt{3471114802157927917746839657246275627496822167060489399554910203}}$\\
\ExactLedgerNext
$xy^{2}z^{2}$ & $U$ & $\frac{\mathtt{1438698269678151371696913059361694896892235546906064110725007289965}}{\mathtt{6626673713210589661153057527470162561584842318933661580968464933}}$\\
$xy^{2}z^{2}$ & $W$ & $\frac{\mathtt{-80096480704524798441566620354536123544154597802973268379738357665}}{\mathtt{946667673315798523021865361067166080226406045561951654424066419}}$\\
$xyz^{3}$ & $U$ & $\frac{\mathtt{3102145252841255952780604187946926413931338657791125868394786111962775}}{\mathtt{220429674396237054488595305593767487448558194897009318829335017531312}}$\\
$xyz^{3}$ & $W$ & $\frac{\mathtt{63124143098168155026060265986978624933421514365737963701974316645275}}{\mathtt{220429674396237054488595305593767487448558194897009318829335017531312}}$\\
$xz^{4}$ & $U$ & $\mathtt{0}$\\
$xz^{4}$ & $W$ & $\mathtt{0}$\\
$y^{5}$ & $U$ & $\frac{\mathtt{186871683168118003488935611555116889880084010}}{\mathtt{32767303919729427758399266748957493087161279}}$\\
$y^{5}$ & $W$ & $\frac{\mathtt{25444070085364078630942054472628302788437558}}{\mathtt{32767303919729427758399266748957493087161279}}$\\
\ExactLedgerNext
$y^{4}z$ & $U$ & $\mathtt{2}$\\
$y^{4}z$ & $W$ & $\mathtt{0}$\\
$y^{3}z^{2}$ & $U$ & $\mathtt{0}$\\
$y^{3}z^{2}$ & $W$ & $\mathtt{0}$\\
$y^{2}z^{3}$ & $U$ & $\mathtt{0}$\\
$y^{2}z^{3}$ & $W$ & $\mathtt{0}$\\
$yz^{4}$ & $U$ & $\mathtt{0}$\\
$yz^{4}$ & $W$ & $\mathtt{0}$\\
$z^{5}$ & $U$ & $\mathtt{0}$\\
$z^{5}$ & $W$ & $\mathtt{0}$\\
\ExactLedgerTail
\end{landscape}
% END INLINED EXACT CPLUS T PARAMETER TABLE

The exact canonical-map identity on $C_+$ is
\begin{equation}\label{eq:t-inverse-identity}
                         G_t+tH_t=0,
\end{equation}
and therefore $t=-G_t/H_t$ away from the common projective base locus.  The
certificate verifies \eqref{eq:t-inverse-identity} in the homogeneous
coordinate ring and checks the base points in the three branch fibres.
Direct substitution gives
$t(P_1)=0$, $t(P_2)=1$, and $t(P_3)=t(P_4)=t(P_5)=\infty$.
Local expansions on the smooth quartic give
\[
 \ord_{P_1}(t)=5,\qquad \ord_{P_2}(t-1)=3,\qquad
 \ord_{P_3}(1/t)=\ord_{P_4}(1/t)=7,\qquad \ord_{P_5}(1/t)=1.
\]
In particular, $P_5$ is an unramified point above the branch value
$\infty$.

\section{Integral coordinates for the genus-two formal group}
\label{app:flynn}

On the divisor-pair chart, write
$\alpha_1=x_1+x_2$ and $\alpha_2=x_1x_2$.  The unnormalized rational
coordinate functions in Flynn's projective embedding include
\[
 a_{15}=1,\quad a_{14}=\alpha_1,\quad a_{13}=\alpha_2,\quad
 a_{12}=\alpha_1^2-2\alpha_2,\quad
 a_{11}=\alpha_1\alpha_2,\quad a_{10}=\alpha_2^2.
\]
These formulas use the representative with $a_{15}=1$; they do not
normalize $a_0$ to $1$.  Near the projective origin, the normalized
coordinates are instead $s_i=a_i/a_0$.  In particular,
$a_0a_{15}=a_5^2$, so $s_{15}=s_5^2$ and $s_{15}$ vanishes at the origin.

The computations use the following integral coordinate block:
\[
 b_i=a_i\quad(0\leq i\leq11),\qquad
 b_{12}=a_{13},\quad b_{13}=a_{14},\quad b_{14}=a_{15},\quad
 b_{15}=a_{12}-2a_{13}.
\]
Its inverse is $a_{12}=b_{15}+2b_{12}$, with all other displayed
coordinates recovered directly.  Hence this is an integral unimodular
change of basis.  Membership in the depth-two origin neighbourhood is
therefore equivalent in the $a$- and $b$-coordinates.  The valuation rows
in \eqref{eq:section7-valuations} are stated in the $b$-coordinates.

\section{The split-prime data at 23 for the rational-parameter theorem}
\label{app:p23-data}

The prime $23$ splits in $K_7$ as the two embeddings $s\mapsto4$ and
$s\mapsto19$.  Both reductions of $C_+$ have $20$ rational points.  Their
common local numerator of the zeta function is
\[
 1-4T+4T^2+56T^3+92T^4-2116T^5+12167T^6,
\]
so
\[
                      \#J_+(\F_{23})=10200,\qquad 23\nmid10200.
\]
The ordered reductions of the five points are
\[
\begin{array}{c|cc}
&s=4&s=19\\ \hline
P_1&(1,21,1)&(1,11,15)\\
P_2&(1,0,22)&(1,8,14)\\
P_3&(1,9,0)&(1,13,0)\\
P_4&(1,1,0)&(1,9,0)\\
P_5&(1,11,21)&(1,6,9).
\end{array}
\]
For the regular differential basis $\eta_x,\eta_y,\eta_z$, the corresponding
evaluation vectors are
\[
\begin{array}{c|cc}
&s=4&s=19\\ \hline
P_1&(3,17,3)&(11,6,4)\\
P_2&(4,0,19)&(5,17,1)\\
P_3&(4,13,0)&(2,3,0)\\
P_4&(12,12,0)&(22,14,0)\\
P_5&(7,8,9)&(10,14,21).
\end{array}
\]
These data independently reproduce the tangent determinants in
\eqref{eq:siksek-dets} once the two annihilator rows are applied.

\small
\begin{thebibliography}{99}

\bibitem{DahmenSiksek}
S.~R. Dahmen and S.~Siksek,
\emph{On the generalized Fermat equation $x^5+y^3=z^7$},
working paper, July 2024,
\url{https://few.vu.nl/~sdn249/GFE357.pdf}.

\bibitem{Putz}
C.~Putz,
\emph{Enumeration of local and global etale algebras applied to generalized
Fermat equations}, PhD thesis, Vrije Universiteit Amsterdam, 2024,
\href{https://doi.org/10.5463/thesis.832}{doi:10.5463/thesis.832}.

\bibitem{PVT}
A.~Pacetti and L.~Villagra Torcomian,
\emph{On the generalized Fermat equation of signature $(5,p,3)$},
arXiv:2512.17845,
\url{https://arxiv.org/abs/2512.17845}.

\bibitem{BreuilDiamond}
C.~Breuil and F.~Diamond,
\emph{Formes modulaires de Hilbert modulo $p$ et valeurs d'extensions entre
caract\`eres galoisiens},
Ann. Sci. \'{E}c. Norm. Sup\'{e}r. (4) \textbf{47} (2014), 905--974.

\bibitem{Raynaud1974}
M.~Raynaud,
\emph{Sch\'emas en groupes de type $(p,\ldots,p)$},
Bull. Soc. Math. France \textbf{102} (1974), 241--280,
\href{https://doi.org/10.24033/bsmf.1779}{doi:10.24033/bsmf.1779}.

\bibitem{DembeleVoight}
L.~Demb\'el\'e and J.~Voight,
\emph{Explicit methods for Hilbert modular forms},
\href{https://jvoight.github.io/articles/hmf-crm-bcn-053024.pdf}{author's version},
Theorem~3.9 and \S4.

\bibitem{Flynn}
E.~V. Flynn,
\emph{The Jacobian and formal group of a curve of genus $2$ over an
arbitrary ground field},
Math. Proc. Cambridge Philos. Soc. \textbf{107} (1990), 425--441,
\href{https://doi.org/10.1017/S0305004100068729}
     {doi:10.1017/S0305004100068729}.

\bibitem{CasselsFlynn}
J.~W.~S. Cassels and E.~V. Flynn,
\emph{Prolegomena to a Middlebrow Arithmetic of Curves of Genus $2$},
London Mathematical Society Lecture Note Series 230,
Cambridge University Press, 1996.

\bibitem{BPS}
N.~Bruin, B.~Poonen, and M.~Stoll,
\emph{Generalized explicit descent and its application to curves of genus
$3$}, Forum Math. Sigma \textbf{4} (2016), e6.

\bibitem{Siksek}
S.~Siksek,
\emph{Explicit Chabauty over number fields},
Algebra Number Theory \textbf{7} (2013), 765--793.

\bibitem{BruinEC}
N.~Bruin,
\emph{Chabauty methods using elliptic curves},
J. Reine Angew. Math. \textbf{562} (2003), 27--49.

\bibitem{BruinStollDescent}
N.~Bruin and M.~Stoll,
\emph{Two-cover descent on hyperelliptic curves},
Math. Comp. \textbf{78} (2009), 2347--2370.

\bibitem{BruinStollSieve}
N.~Bruin and M.~Stoll,
\emph{The Mordell--Weil sieve: proving non-existence of rational points on
curves},
LMS J. Comput. Math. \textbf{13} (2010), 272--306.

\bibitem{Magma}
W.~Bosma, J.~Cannon, and C.~Playoust,
\emph{The Magma algebra system. I. The user language},
J. Symbolic Comput. \textbf{24} (1997), 235--265.

\bibitem{PARI}
The PARI Group,
\emph{PARI/GP},
\url{https://pari.math.u-bordeaux.fr/}.

\bibitem{Sage}
The Sage Developers,
\emph{SageMath, the Sage Mathematics Software System},
\url{https://www.sagemath.org/}.

\end{thebibliography}
\normalsize

\end{document}
